\section*{Chapitre \ref{ch:g2:what-animals-eat} --- Ce que mangent les animaux}

\begin{solution}{exo:g2:what-animals-eat:1}
\omterm{def:g2:what-animals-eat:kinds}{Herbivore}~; \omterm{def:g2:what-animals-eat:kinds}{carnivore}~; \omterm{def:g2:what-animals-eat:kinds}{omnivore}.
\end{solution}

\begin{solution}{exo:g2:what-animals-eat:2}
(a) \omterm{def:g2:what-animals-eat:kinds}{Herbivore}. (b) \omterm{def:g2:what-animals-eat:kinds}{Carnivore}. (c) \omterm{def:g2:what-animals-eat:kinds}{Omnivore}. (d) \omterm{def:g2:what-animals-eat:kinds}{Herbivore}. (e)
\omterm{def:g2:what-animals-eat:kinds}{Omnivore}.
\end{solution}

\begin{solution}{exo:g2:what-animals-eat:3}
Les broyeuses plates --- les \omterm{def:g1:plants-around-us:plant}{plantes} comme l'herbe sont coriaces et
doivent être broyées longtemps avant d'être avalées~; il n'y a
aucune proie à saisir.
\end{solution}

\begin{solution}{exo:g2:what-animals-eat:4}
Court et épais~: casser les \omterm{def:g2:from-seed-to-plant:seed}{graines}. Courbé et crochu~: déchirer la
viande. Long et fin~: aller chercher les insectes dans l'écorce et
les fentes.
\end{solution}

\begin{solution}{exo:g2:what-animals-eat:5}
De minuscules pucerons --- la coccinelle en mange des quantités, ce
qui la fait adorer des jardiniers.
\end{solution}

\begin{solution}{exo:g2:what-animals-eat:6}
L'herbe nourrit peu à chaque bouchée, alors la vache doit brouter
pendant des heures pour en avoir assez~; la viande est riche, alors
un bon repas peut tenir le renard un long moment.
\end{solution}

\begin{solution}{exo:g2:what-animals-eat:7}
Des vers et des insectes en été~; des baies en hiver, quand le sol
est trop dur pour les vers. Un large menu permet à un \omterm{def:g2:what-animals-eat:kinds}{omnivore} de
changer d'aliment quand l'un vient à manquer --- une grande aide
dans les saisons difficiles.
\end{solution}

\begin{solution}{exo:g2:what-animals-eat:8}
Réponse libre. En général~: moineaux et pinsons (becs courts et
épais) prennent des \omterm{def:g2:from-seed-to-plant:seed}{graines}~; les mésanges attrapent insectes et
\omterm{def:g2:from-seed-to-plant:seed}{graines}~; les merles fouillent pour des vers et prennent des baies.
L'accord bec--menu de la proposition se vérifie bien.
\end{solution}

\begin{solution}{exo:g2:what-animals-eat:9}
Le menu d'un \omterm{def:g2:what-animals-eat:kinds}{carnivore}~: d'autres \omterm{def:g1:animals-around-us:animal}{animaux}. Les étapes 1 et 2 --- les
dents (des pointes qui saisissent) et les outils de chasse (griffes,
\omterm{def:g1:the-five-senses:sense}{yeux} vers l'avant)~; l'étape 4 lui a donné son nom.
\end{solution}

\begin{solution}{exo:g2:what-animals-eat:10}
Le renard mange des lapins, et les lapins mangent de l'herbe. Pas
d'herbe, pas de lapins~; pas de lapins, des renards affamés. Le
renard dépend de l'herbe \emph{à travers} ses proies --- dans la
nature, les repas forment des chaînes.
\end{solution}
